% !TEX program = lualatex
\documentclass[12pt,a4paper]{article}

% ========= 语言包 =========
\usepackage[english]{babel}
\babelprovide[import, main]{chinese-simplified}
\babelprovide[import, onchar=ids fonts]{english}
\usepackage{fontspec}
\defaultfontfeatures{Ligatures=TeX}

% ========= 中文字体 (Noto CJK SC) =========
\babelfont[chinese-simplified]{rm}[Renderer=Harfbuzz]{Noto Serif CJK SC}
\babelfont[chinese-simplified]{sf}[Renderer=Harfbuzz]{Noto Sans CJK SC}
\babelfont[chinese-simplified]{tt}[Renderer=Harfbuzz]{Noto Sans Mono CJK SC}

% ========= 英文字体 (标准系统字体) =========
\babelfont[english]{rm}{Times New Roman}
\babelfont[english]{sf}{Arial}
\babelfont[english]{tt}{Courier New}

% ========= 宏包 =========
\usepackage{amsmath,amssymb,amsfonts,mathtools}
\usepackage{geometry}
\geometry{left=1.0cm,right=1.0cm,top=1.0cm,bottom=3.4cm}
\usepackage{booktabs}
\usepackage{hyperref}
\usepackage{url}
\usepackage{xcolor}
\usepackage{titlesec}
\usepackage{fancyhdr}
\usepackage{microtype}
\usepackage{array}
\usepackage{tikz}
\usetikzlibrary{arrows.meta, positioning, calc, shapes.geometric, decorations.pathreplacing}

% ========= YUCT 颜色 =========
\definecolor{skyblue}{RGB}{0,102,204}
\definecolor{darkblue}{RGB}{0,0,139}
\definecolor{gold}{RGB}{255,215,0}

% ========= 章节格式 =========
\titleformat{\section}{\LARGE\bfseries\color{darkblue}}{\thesection}{1em}{}
\titleformat{\subsection}{\Large\bfseries\color{darkblue}}{\thesubsection}{1em}{}

% ========= 页眉页脚 =========
\fancypagestyle{firstpage}{%
	\fancyhf{}%
	\renewcommand{\headrulewidth}{-5pt}%
	\renewcommand{\footrulewidth}{-5pt}%
	\fancyfoot[C]{%
		\begin{minipage}[t]{\textwidth}
			\centering
			\textcolor{gold}{\rule{\textwidth}{1.6pt}}\\[5pt]
			{\small \textsuperscript{1}Yakushev Research, \url{https://yuct.org/}} \hfill
			{\small YPSDC 协议: \url{https://ypsdc.com/}}\\[-2pt]
			\textcolor{gold}{\rule{\textwidth}{1.6pt}}\\[5pt]
			\textbf{雅库舍夫统一协调理论 2026} \hfill \thepage\\[10pt]
		\end{minipage}%
	}%
}

\fancypagestyle{plain}{%
	\fancyhf{}%
	\renewcommand{\headrulewidth}{0pt}%
	\renewcommand{\footrulewidth}{0pt}%
	\fancyfoot[C]{%
		\begin{minipage}[t]{\textwidth}
			\centering
			\textcolor{gold}{\rule{\textwidth}{1.6pt}}\\[4pt]
			\textbf{雅库舍夫统一协调理论 2026} \hfill \thepage\\[4pt]
		\end{minipage}%
	}%
}

\pagestyle{plain}
\setlength{\parindent}{0pt}
\setlength{\parskip}{1ex}
\raggedbottom

\hypersetup{
	colorlinks=true,
	linkcolor=blue,
	citecolor=blue,
	urlcolor=blue,
}

% ========= 页眉 (YUCT 标志) =========
\newcommand{\makeheader}{%
	\noindent
	\begin{minipage}{0.17\textwidth}
		\raggedright
		{\sffamily\bfseries\color{skyblue}\fontsize{46}{46}\selectfont YUCT}%
	\end{minipage}%
	\hspace{30pt}
	\begin{minipage}{0.32\textwidth}
		\raggedright
		{\sffamily\fontsize{11}{13}\selectfont 雅库舍夫\\ 统一\\ 协调理论}%
	\end{minipage}%
	\hfill
	\begin{minipage}{0.38\textwidth}
		\raggedleft
		{\sffamily\bfseries 假设 / 路线图}\\ 
		{\sffamily 版本 2.0 / 2026年6月}\\ 
	\end{minipage}
	\par\vspace{5pt}
	\noindent\textcolor{gold}{\rule{\textwidth}{1.6pt}}
	\par\vspace{0pt}
}

\begin{document}
	\thispagestyle{firstpage}
	\makeheader
	
	\vspace{0cm}
	{\LARGE\bfseries 强CP问题：\\[0.1cm] YUCT框架内的一个具体代数假设\\[0.2cm] \Large 理论与实验验证路线图 \par}
	\vspace{0.2cm}
	
	{\large 阿列克谢·V·雅库舍夫\textsuperscript{1}} \\
	\vspace{0.0cm}
	
	{\color{darkblue}\textbf{\LARGE 摘要}} \par
	{\large
		\noindent
		强CP问题——为什么QCD中破坏CP的\(\theta\)参数小于\(10^{-10}\)——是雅库舍夫统一协调理论(YUCT)中标准模型的最后一个未解谜题。在本文中，我们提出一个具体的代数假设，自然产生所需的抑制。我们假设\(\theta\)参数的有效协调深度为
		\[
		L_\theta = \frac{1}{2}\,L_W \;+\; \frac{S_{\mathrm{odd}}}{S_{\mathrm{even}}},
		\]
		其中\(L_W = 96\)是\(W\)玻色子的深度，\(S_{\mathrm{odd}}/S_{\mathrm{even}} = 3/2\)是理论的基本常数。该公式不含自由参数，给出\(L_\theta = 49.5\)，对应\(\theta \sim (2/3)^{49.5} \approx 2.3 \times 10^{-10}\)。我们讨论该表达式的物理动机，将其嵌入YUCT的更广泛代数框架，并概述一个实验测试计划——主要通过中子电偶极矩(nEDM)——可以证伪该假设。该提议作为一个清晰、数学精确的猜想提出，有待未来数据证伪或确认。
	}
	
	\vspace{0.1cm}
	\noindent
	{\color{darkblue}\textbf{关键词:}} YUCT, 强CP问题, 代数假设, 协调深度, 中子电偶极矩, 开放研究计划。
	\vspace{0.4cm}
	
	\clearpage
	\tableofcontents
	\newpage
	
	\section{引言}
	\label{sec:intro}
	
	雅库舍夫统一协调理论(YUCT)已经为等级问题、宇宙常数问题、费米子质量模式以及基础物理学的若干其他重大谜题提供了定量的、无参数的解决方案~\cite{Y1,Y3,Y4}。唯一的例外一直是强CP问题：量子色动力学(QCD)中破坏CP的\(\theta\)参数的人为小量，其实验上限为\(|\theta| \lesssim 10^{-10}\)~\cite{nEDM2020}。
	
	在YUCT内部，早期试图为\(\theta\)分配一个协调深度的尝试——通过一个临时的整数\(k_\theta\)——是不令人满意的，因为它们退化为拟合而非结构性预测。在本笔记中，我们提出一种不同的方法：我们利用QCD真空的有序(费米子)与混沌(拓扑)扇区之间的相互作用——YUCT凭借其对有序和混沌的统一代数描述，在量化这种相互作用方面具有独特优势~\cite{Y2,Feig}。
	
	我们不声称从19维拉格朗日量进行了严格推导；这样的推导是一个长期目标。相反，我们陈述一个尖锐的代数猜想，它自然地源于YUCT中已经存在的数字，不含可调参数，并对中子电偶极矩(nEDM)和轴子质量做出了可证伪的预测。该猜想作为一个定义明确的目标，供未来的理论和实验工作研究。
	
	\section{YUCT中的有序与混沌常数}
	\label{sec:oc}
	
	\subsection{有序}
	普遍误差律
	\[
	\varepsilon = \kappa_c\,\alpha\,K_{\mathrm{eff}}^{-\beta},\qquad \beta = \frac{2}{3},\quad \kappa_c = \frac{1}{3},
	\]
	与协调的层级结构一起，导致封闭的代数循环
	\[
	S_{\mathrm{odd}} = \frac{6}{5} = 1.2,\quad S_{\mathrm{even}} = \frac{4}{5} = 0.8,\quad 
	\frac{S_{\mathrm{odd}}}{S_{\mathrm{even}}} = \frac{3}{2} = \frac{1}{\beta}.
	\]
	所有物理质量都用普朗克质量和协调深度\(L\)表示：
	\[
	m \approx M_{\mathrm{Pl}}\left(\frac{2}{3}\right)^L.
	\]
	特别地，\(W\)玻色子的质量对应于\(L_W = 96\)，即标准模型中韦尔费米子自由度的总数(包括右手中微子)~\cite{Y3}。
	
	\subsection{混沌}
	普遍费根鲍姆常数
	\[
	\delta \approx 4.6692,\qquad \alpha \approx 2.5029,
	\]
	描述了通向混沌的倍周期道路。在YUCT中，它们通过黄金比例\(\varphi = (1+\sqrt{5})/2\)和圆周率\(\pi\)与有序常数代数地联系起来：
	\[
	2\alpha^2 \approx \pi\delta,\qquad
	\pi \approx S_{\mathrm{odd}}\varphi^2 \quad \Longrightarrow \quad 
	2\alpha^2 \approx (S_{\mathrm{odd}}\varphi^2)\,\delta .
	\]
	这种联系表明费根鲍姆常数并非外在于YUCT；它们属于同一扩展代数系统~\cite{Y2}。
	
	\subsection{量子真空作为有序-混沌共轭}
	QCD真空是有序方面(夸克凝聚、手征对称性破缺)和混沌方面(瞬子涨落、拓扑磁化率)共存的舞台。\(\theta\)参数与拓扑荷密度\(G\widetilde{G}\)耦合，后者从两个扇区都获得贡献。因此，很自然地期望有效深度\(L_\theta\)是表征有序和混沌分量的深度的函数。
	
	在先前的工作中，我们引入了初步深度\(L_{\mathrm{ord}}\)(费米子扇区)和\(L_{\mathrm{ch}}\)(拓扑/混沌扇区)，但它们的简单线性组合并未重现所需的抑制\(|\theta| \sim 10^{-10}\)。这促使我们寻找更具结构性的组合。
	
	\section{关于\(L_\theta\)的代数假设}
	\label{sec:hypothesis}
	
	\subsection{假设的陈述}
	
	我们提出\(\theta\)参数的有效协调深度由简单公式给出
	\begin{equation}
		\boxed{L_\theta = \frac{1}{2}\,L_W \;+\; \frac{S_{\mathrm{odd}}}{S_{\mathrm{even}}}},
		\label{eq:Ltheta}
	\end{equation}
	其中
	\begin{itemize}
		\item \(L_W = 96\)是\(W\)玻色子的深度(电弱尺度的基线)；
		\item \(S_{\mathrm{odd}}/S_{\mathrm{even}} = 3/2 = 1.5\)是支配代间质量因子的普遍比率，出现在初级代数循环中。
	\end{itemize}
	
	代入这些数字，
	\[
	L_\theta = \frac{96}{2} + 1.5 = 48 + 1.5 = 49.5.
	\]
	
	相应的\(\theta\)参数，最多乘以一个来自瞬子计算的阶为1的因子，为
	\[
	\theta \sim \left(\frac{2}{3}\right)^{L_\theta} = \left(\frac{2}{3}\right)^{49.5} \approx 2.3 \times 10^{-10}.
	\]
	这正好落在实验允许的窗口\(|\theta| \lesssim 10^{-10}\)内~\cite{nEDM2020}。
	
	\subsection{物理解释}
	
	为什么会出现这种特定的组合？
	\begin{enumerate}
		\item \textbf{因子\(1/2\)}：电弱尺度由真空期望值\(v \approx 246\) GeV设定，通过几何因子\(\xi_W \approx 0.53\)从基础深度96获得。深度48对应于尺度\(\sqrt{v M_{\mathrm{Pl}}}\)——弱尺度与普朗克尺度的几何平均——这正是费米子与玻色子对真空能量的贡献预期交叉的尺度。在YUCT中，已知该尺度在希格斯凝聚的协调动力学中起着特殊作用(附录\(\Lambda\))。
		\item \textbf{偏移量\(S_{\mathrm{odd}}/S_{\mathrm{even}} = 3/2\)}：比率\(3/2 = 1/\beta\)量化了任何层级协调系统中有序(奇数层)与无序(偶数层)关联之间的不可约不对称性。它出现在\(L_\theta\)中表明，CP破坏项源于真空中以有序为主和以混沌为主的扇区之间一个小的、不可避免的不平衡。这个偏移量是加性的而非乘性的，反映了协调网络分立的、分层结构。
	\end{enumerate}
	
	因此，方程~(\ref{eq:Ltheta})可以解读为：\(\theta\)的深度等于弱深度的一半加上普遍的有序-混沌偏移量。它不含自由参数，只使用了YUCT中已经独立确定的数字。
	
	\subsection{直接后果}
	
	\paragraph{中子电偶极矩。}
	nEDM通过\(d_n \approx 3.6 \times 10^{-16}\,\theta\; e\!\cdot\!\mathrm{cm}\)与\(\theta\)相关(含强子不确定性因子)~\cite{nEDM2020}。取\(\theta \approx 2.3 \times 10^{-10}\)给出
	\[
	d_n \approx 8.3 \times 10^{-26}\; e\!\cdot\!\mathrm{cm}.
	\]
	当前的实验界限为\(d_n < 1.8 \times 10^{-26}\; e\!\cdot\!\mathrm{cm}\)(90\%置信度)。下一代的PSI、SNS和TUCAN实验目标灵敏度为\(10^{-28}\; e\!\cdot\!\mathrm{cm}\)量级。该水平上的零结果将排除该假设；而正检测将是戏剧性的确认。
	
	\paragraph{轴子物理。}
	如果强CP问题通过Peccei-Quinn机制解决，则轴子质量\(m_a\)与\(\theta\)及轴子衰变常数\(f_a\)相关。方程~(\ref{eq:Ltheta})通过相同的协调深度逻辑确定了\(f_a\)(从而\(m_a\))。这为ADMX和其他轴子搜索提供了互补的实验目标。
	
	\section{验证路线图}
	\label{sec:roadmap}
	
	该假设在数学上是精确且可证伪的。我们提出以下计划：
	
	\subsection{理论任务}
	\begin{enumerate}
		\item \textbf{从19维拉格朗日量推导。}
		必须将YUCT的19维母理论的有效作用量约化为标准的QCD\(G\widetilde{G}\)项。目标是表明，当识别出正确的真空态时，约化中的系数自然产生深度\(L_\theta = 49.5\)。这是一个具有挑战性但定义明确的数学问题。
		\item \textbf{格点规范理论。}
		现有的格点QCD计算可以计算拓扑磁化率和真空能量对\(\theta\)的依赖关系。受YUCT启发的分析应测试方程~(\ref{eq:Ltheta})的代数结构是否反映在格点结果随耦合常数的标度行为中，或反映在拓扑荷簇的分布中。
		\item \textbf{精化nEDM预测。}
		将\(\theta\)与\(d_n\)联系起来的强子矩阵元已从手征微扰论和格点QCD中得知。对这些矩阵元的专门计算，结合YUCT给出的\(\theta\)值，将产生一个尖锐的预测，可直接与下一代nEDM实验的界限比较。
	\end{enumerate}
	
	\subsection{实验任务}
	\begin{enumerate}
		\item \textbf{nEDM实验(PSI, SNS, TUCAN)。}
		测量的nEDM高于\(10^{-26}\;e\!\cdot\!\mathrm{cm}\)将与假设一致；低于\(10^{-28}\;e\!\cdot\!\mathrm{cm}\)的界限将排除它。
		\item \textbf{轴子搜索(ADMX, MADMAX, IAXO)。}
		如果Peccei-Quinn机制得以实现，那么一旦从\(L_\theta\)评估出\(f_a\)，YUCT预测的轴子质量应是可计算的。在相应质量窗口中的专门搜索将检验该假设。
		\item \textbf{宇宙学印记。}
		如果有序-混沌干涉是基本的，它也可能影响早期宇宙的相变，在宇宙微波背景或暗物质分布中留下特定印记。这些可以通过Planck、Euclid和未来的CMB-S4数据加以约束。
	\end{enumerate}
	
	\section{合作邀请}
	\label{sec:invite}
	
	YUCT框架，包括其完整的拉格朗日量、所有附录以及协调深度的广泛数据库，已在Zenodo上开放获取~\cite{Y1,Y2}。我们邀请物理学家、数学家和计算科学家参与推导、检验或证伪本文提出的假设。作者(阿列克谢·雅库舍夫，\href{mailto:alexey@yakushev.eu}{alexey@yakushev.eu})欢迎通信、建议和合作提案。
	
	\section{结论}
	\label{sec:conclusion}
	
	我们在雅库舍夫统一协调理论中提出了一个具体的、无参数的代数公式，用于QCD\(\theta\)参数的有效深度：
	\[
	L_\theta = \frac{1}{2}L_W + \frac{S_{\mathrm{odd}}}{S_{\mathrm{even}}} = 49.5,
	\]
	给出\(\theta \approx 2.3 \times 10^{-10}\)。这一假设在数学上简单，在物理上具有动机，并且可通过即将进行的nEDM和轴子实验直接证伪。它取代了依赖任意整数指派的早期尝试，并为在YUCT内完全解决强CP问题开辟了清晰的路径。该假设是否能经受住实验的检验还有待观察；其价值在于提供一个精确的、可检验的猜想，以指导未来的研究。
	
	\vspace{1cm}
	\noindent\textbf{日期:} 2026年6月\\
	\textbf{版本:} 2.0\\
	\textbf{状态:} 具有特定代数预测的假设；开放进行实验检验和理论推导。
	
	\newpage
	\begin{thebibliography}{99}
		\bibitem{Y1} A.V. Yakushev, \textit{Yakushev Unified Coordination Theory (YUCT)}, Zenodo, 2026. DOI:10.5281/zenodo.18444599.
		\bibitem{Y2} A.V. Yakushev, Appendix AF: The Algebraic Framework of Reality in YUCT, 2026.
		\bibitem{Y3} A.V. Yakushev, Appendix \(\Lambda\): Resolution of the Hierarchy and Cosmological Constant Problems, 2026.
		\bibitem{Y4} A.V. Yakushev, Appendix J: Complete Derivation of Standard Model Flavor Parameters, 2026.
		\bibitem{Feig} M.J. Feigenbaum, \textit{J. Stat. Phys.} \textbf{19}, 25 (1978).
		\bibitem{nEDM2020} C. Abel et al. (nEDM Collaboration), \textit{Phys. Rev. Lett.} \textbf{124}, 081803 (2020).
	\end{thebibliography}
	
\end{document}